% PASS20260921 E8 matter-81 complete-frame x external-qutrit tensor carrier.
\subsection*{The $81$-root matter shell is $27$ complete frames times one external qutrit}

The exact $E_8\to E_6\times A_2$ coordinate decomposition gives
\[
81=27_{E_6}\times3_{A_2}
\]
root by root. Independently, the two-qutrit frame theorem identifies the cubic-surface $27$ with the $27$ complete tensor-factorisation frames.

In one frozen matter chart, an anchored complement-Schl\"afli isomorphism assigns every root a unique label
\[
\boxed{(F,a),\qquad F\in\{0,\ldots,26\},\quad a\in\mathbb F_3.}
\]
The cubic law factorizes exactly. The $27$ projected weights have $45$ zero-sum tritangents. For each tritangent all $3^3$ phase assignments are tested; exactly six lift to zero-sum $E_8$ triples, namely the six permutations of the three distinct phases. Hence
\[
\boxed{270=45\cdot3!}
\]
zero-sum matter triples occur, and every one of the $81$ roots lies in
\[
\boxed{10=5\cdot2}
\]
of them. Cubic root addition therefore enforces the external-qutrit selection rule: one of each phase.

\paragraph{Scope firewall.}
This is an object-set/root-incidence theorem in a fixed coordinate gauge, not a Hilbert-space factorization of the full adjoint, a Yukawa coefficient, or a vacuum construction.
